\section*{Author exact-clique model and SOS conventions}
The SOS proof system, also known as the Positivstellensatz system [14], works with polynomials over $\mathbb R$. Given polynomial equalities (axioms) $f_1(\mathbf x)=0,\ldots,f_k(\mathbf x)=0$ on $\mathbf x=(x_1,\ldots,x_n)$, a proof (a refutation of the existence of a solution) is
\[-1=\sum_{i=1}^kf_iq_i+\sum_jr_j^2\quad\text{in }\mathbb R[x_1,\ldots,x_n],\]
where $q_1,\ldots,q_m$ and $r_1,\ldots$ are arbitrary polynomials on $x_1,\ldots,x_n$ over $\mathbb R$. Under certain conditions, in particular when all variables are boolean ($x_i^2=x_i$), such an refutation always exists if the axioms are contradictory. The degree-$d$ SOS proof system adds the limitation
\[\max_{i,j}\{\deg(f_i)+\deg(q_i),2\deg(r_j)\}\le d.\]
\paragraph{Definition 1.1.} Given an $n$-vertex simple graph $G$ and a number $\omega$, the Clique Problem for degree-$d$ SOS proof system has the following axioms:
\begin{align*}(\text{Boolean})\quad x_i^2&=x_i&&\forall i\in[n],\\(\text{Clique})\quad x_ix_j&=0&&\forall\{i,j\}\text{ non-edge},\\(\text{Size})\quad x_1+\cdots+x_n&=\omega.&&\tag{1.1}\end{align*}
To confirm no $\omega$-clique exists is to give an SOS refutation of the above. The SOS system has duality: to show degree lower bound it suffices to consider pseudo-expectation and the resulting moment matrix.\footnote{The author uses the name for simplicity. More cautiously, it should be called the pseudo-moment matrix.} With boolean variables this can be demonstrated on multi-linear polynomials. Let $X^{\le a}=\{x_S:S\subseteq[n],|S|\le a\}$ for any $a$.
\paragraph{Definition 1.2.} A degree-$d$ pseudo-expectation for the Clique Problem on $G$ is a map $\widetilde{\mathbb E}:X^d\to\mathbb R$ satisfying the following four constraints when extended by $\mathbb R$-linearity:
\begin{align*}(\text{Default})\quad\widetilde{\mathbb E}x_\emptyset&=1,\tag{1.2}\\(\text{Clique})\quad\widetilde{\mathbb E}x_S&=0&&\forall S:|S|\le d,\ G|_S\text{ non-clique},\tag{1.3}\\(\text{Size})\quad\widetilde{\mathbb E}((x_1+\cdots+x_n)x_S)&=\omega\widetilde{\mathbb E}x_S&&\forall S:|S|\le d-1.\tag{1.4}\end{align*}
In (1.4), $x_Ax_B:=x_{A\cup B}$. Define the moment matrix $M$ to be the $\binom{[n]}{\le d/2}\times\binom{[n]}{\le d/2}$ matrix with expression $M(A,B)=\widetilde{\mathbb E}x_{A\cup B}$ for all $|A|,|B|\le d/2$.\footnote{$d$ is always assumed to be even.} The last constraint is:\[M\text{ is positive semi-definite}.\tag{1.5}\]
If a degree-$d$ pseudo-expectation exists then there is no degree-$d$ SOS refutation. Exact Clique is Definition1.1. The Non-Exact Clique relaxation replaces the Size Constraints by one weaker inequality $\widetilde{\mathbb E}(x_1+\cdots+x_n)\ge\omega$. There is no planted clique in the problem's formulation; traditionally it is still called the planted clique problem due to the algorithmic motivation.
